\documentclass[11pt]{article}
\usepackage{preamble}
\setlength{\parskip}{4pt}

\begin{document}

\daybanner{AP Statistics \textperiodcentered\ Topic 7.3 (CED 4.3) \textperiodcentered\ B: 2027-01-15 \textperiodcentered\ E: 2027-03-03 \textperiodcentered\ TEACHER KEY}{One Interval, Two Cutoffs}

\sectionbanner{CED TETHER}

{\small
\begin{itemize}[leftmargin=1.5em,itemsep=2pt,topsep=2pt]
  \item \textbf{Skill 4.B} --- Interpret statistical calculations and findings to assign meaning or assess a claim.
  \item \textbf{Skill 4.D} --- Justify a claim based on a confidence interval.
  \item \textbf{Skill 4.A} --- Make an appropriate claim or draw an appropriate conclusion.
  \item \textbf{EU UNC-4} --- An interval of values should be used to estimate parameters, in order to account for uncertainty.
  \item \textbf{LO UNC-4.S} --- Interpret a confidence interval for a population mean, including the mean difference between values in matched pairs. [Skill 4.B]
  \item \textbf{EK UNC-4.S.1} --- A confidence interval for a population mean either contains the population mean or it does not, because each interval is based on data from a random sample, which varies from sample to sample.
  \item \textbf{EK UNC-4.S.2} --- We are C\% confident that the confidence interval for a population mean captures the population mean.
  \item \textbf{EK UNC-4.S.3} --- An interpretation of a confidence interval for a population mean includes a reference to the sample taken and details about the population it represents.
  \item \textbf{LO UNC-4.T} --- Justify a claim based on a confidence interval for a population mean, including the mean difference between values in matched pairs. [Skill 4.D]
  \item \textbf{EK UNC-4.T.1} --- A confidence interval for a population mean provides an interval of values that may provide sufficient evidence to support a particular claim in context.
  \item \textbf{LO UNC-4.U} --- Identify the relationships between sample size, width of a confidence interval, confidence level, and margin of error for a population mean. [Skill 4.A]
  \item \textbf{EK UNC-4.U.1} --- When all other things remain the same, the width of a confidence interval for a population mean tends to decrease as the sample size increases.
  \item \textbf{EK UNC-4.U.2} --- For a single mean, the width of the interval is proportional to 1/$\surd$n.
  \item \textbf{EK UNC-4.U.3} --- For a given sample, the width of the confidence interval for a population mean increases as the confidence level increases.
\end{itemize}
}
\textbf{Essential question:} How can one confidence interval support a bounded claim while leaving a stronger claim unresolved?\par
\textbf{DOK-3 skill:} 4.D \textperiodcentered\ \textbf{FRQ pattern:} {\small\texttt{mean-\allowbreak{}interval-\allowbreak{}two-\allowbreak{}cutoffs-\allowbreak{}and-\allowbreak{}scope}}\par
\textbf{DOK rationale:} Part (c) coordinates both endpoints, two cutoffs, random-sample scope, and the mean-versus-individual distinction to adjudicate two interpretations and trace an advertising consequence.\par

\frameworkphaseheader{Do Now (first take) $\rightarrow$ Explore (video + follow-along) $\rightarrow$ Exit (finish + turn in)}{1 $\rightarrow$ 3}{45}{%
\begin{itemize}[leftmargin=1.5em,itemsep=2pt,topsep=2pt]
  \item Board slide up at the bell. Ask which interval feature evaluates each cutoff.
  \item Begin the 7.3 video follow-along at minute 5.
  \item Return at minute 33. Require separate decisions for the two cutoffs and the individual-time claim.
  \item Collect at the door. Score part (c) only, E/P/I.
\end{itemize}
}{%
\begin{itemize}[leftmargin=1.5em,itemsep=2pt,topsep=2pt]
  \item Commit to one approach and cite an endpoint or the parameter.
  \item Complete the follow-along, finish (a)--(c), revise, and turn in.
\end{itemize}
}{%
\begin{itemize}[leftmargin=1.5em,itemsep=2pt,topsep=2pt]
  \item Why does 44.877 settle the 45-second claim but not the 42-second claim?
  \item Does this interval describe a population mean or 95 percent of individual transactions?
  \item Which population can the random sample represent?
\end{itemize}
}{Solo teacher. Test each cutoff against the interval, then separate the mean from individual times.}

\begin{callout}{\IconWarn\ WATCH FOR}{calloutred}
\begin{itemize}[leftmargin=1.5em,itemsep=2pt,topsep=2pt]
  \item Using only the lower endpoint to support a less-than claim.
  \item Treating 95 percent confidence as a probability that the fixed population mean is in this realized interval.
  \item Claiming that 95 percent of individual setup times fall between the confidence-interval endpoints.
\end{itemize}
\end{callout}

\sectionbanner{SCORING PART (c) --- E / P / I}

\textbf{Expected elements}
\begin{itemize}[leftmargin=1.5em,itemsep=2pt,topsep=2pt]
  \item \textbf{chooses-warranted-approach} (required) --- Selects the full-interval, population-mean interpretation
  \item \textbf{separates-two-cutoffs} (required) --- Uses 44.877 below 45 and 42 inside the interval
  \item \textbf{uses-whole-interval} (required) --- Applies the entire-interval rule for a less-than claim
  \item \textbf{distinguishes-mean-from-individuals} (required) --- Separates the population mean from individual times
  \item \textbf{traces-advertising-consequence} (required) --- Names both unsupported advertising claims
\end{itemize}

{\small\renewcommand{\arraystretch}{1.25}
\noindent\begin{tabular}{|p{0.5in}|p{5.9in}|}
\hline
\textbf{E} & Uses both cutoff locations, the whole-interval rule, the correct parameter, and both advertising consequences. \\ \hline
\textbf{P} & Evaluates both cutoffs but incompletely explains interval logic or the mean-versus-individual distinction. \\ \hline
\textbf{I} & Uses only values below a cutoff or treats the interval as covering 95 percent of individual times. \\ \hline
\end{tabular}}

\clearpage
\focusbanner{TODAY'S PROBLEM --- annotated}

Suppose a library wants the mean setup time for all 800 kiosk transactions in a month below 45 seconds and hopes to advertise a mean below 42 seconds. An SRS of 25 has the shown summary. A 95\% one-sample $t$-interval is $(39.923,44.877)$ seconds.\par\smallskip \textbf{Mina:} ``I would compare each cutoff with the full interval and keep the conclusion about the population mean.''\par \textbf{Gabe:} ``Because the interval contains times below both cutoffs, I think it supports both claims; I also read 95\% as the share of transaction times inside it.''\par
\begin{center}
\textbf{Kiosk sample}\par\smallskip
\begin{tabular}{|l|c|}
\hline
\textbf{Summary} & \textbf{Value} \\ \hline
\textbf{Population} & 800 \\ \hline
\textbf{$n$} & 25 \\ \hline
\textbf{Mean (s)} & 42.4 \\ \hline
\textbf{SD (s)} & 6.0 \\ \hline
\textbf{Shape} & Symmetric; no outliers \\ \hline
\end{tabular}
\end{center}
\begin{firsttakebox}{FIRST TAKE (5 min)}
Which approach is more defensible? Cite an endpoint or identify the quantity being estimated.\par
\answer{Not graded. Students should identify the parameter and an interval feature before deciding.}
\end{firsttakebox}

\textit{Rules callout ``LET THE ENTIRE INTERVAL GOVERN THE CLAIM (keep on the page)'' is printed on the student sheet.}\par\medskip

\dokbadge{1}~\textbf{(a)} Define the population mean and locate 42 and 45 relative to the interval.\par
\answer{Let $\mu$ be the mean setup time for all 800 kiosk transactions that month. The value 42 is inside $(39.923,44.877)$; 45 is above the upper endpoint.}

\dokbadge{2}~\textbf{(b)} Verify random, 10 percent, and shape conditions. Reproduce the standard error, margin of error, and interval; interpret it in context.\par
\answer{The SRS satisfies $25\leq0.10(800)=80$. With a symmetric, outlier-free sample and $n<30$, the $t$ shape condition is reasonable. $SE=6/\sqrt{25}=1.2$ and $ME=2.063899(1.2)=2.4766788$, so the interval is $42.4\pm2.4766788=(39.9233,44.8767)$. We are 95\% confident that the mean for all 800 transactions is between about 39.923 and 44.877 seconds.}

\dokbadge{3}~\textbf{(c)} Adjudicate the approaches. Use both endpoints to evaluate the two less-than claims, explain why the cutoff locations matter, distinguish the population mean from individual times, and trace what the rejected interpretation could lead the library to advertise.\par
\answer{Mina is warranted. The whole interval is below 45, so it supports $\mu<45$. Because 42 is inside, plausible means lie on both sides and $\mu<42$ is not established. A less-than claim needs the entire interval below its cutoff, not merely some values. The interval estimates the population mean represented by the SRS; it does not contain 95\% of individual times. The rejected interpretation could produce unsupported advertisements about both a mean below 42 and individual transaction times.}

\begin{summaryexitbox}{EXIT}
One thing the video changed about my first take: \hrulefill
\end{summaryexitbox}

\end{document}
